\documentclass[11pt,a4paper]{article}
\usepackage[T1]{fontenc}
\usepackage{lmodern}
\usepackage[margin=24mm]{geometry}
\usepackage{amsmath,amssymb}
\usepackage[hidelinks]{hyperref}
\usepackage{microtype}
\newcommand{\R}{\mathbb R}
\newcommand{\g}{\mathfrak g}
\title{Isotropy types are not counted by coadjoint orbits}
\author{Submission for TLMC problem 00000006034}
\date{24 October 2026}
\begin{document}
\maketitle

\paragraph{Scope and result.}
The final clause of the source identifies the number of isotropy types in
nonfree reduction with the number of coadjoint orbits. In the standard
Hamiltonian setting this is false: the sign action of the compact Lie
group $C_2$ on the connected symplectic plane has two occurring isotropy
types in its zero-level quotient, whereas its coadjoint action has one
orbit. We prove this necessary-clause counterexample, without assigning
values to the source's other unspecified constants or quantization data.
The source imposes neither connectedness of the group nor properness of
the moment map. Our group is disconnected, and our moment map is not
proper. The claim here is expressly confined to the printed scope.

\paragraph{An actual smooth symplectic action.}
Let $G=C_2=\{1,s\}$, with $s^2=1$, and give it the discrete topology and
the zero-dimensional real smooth manifold structure. Let $V=\R^2$ with
its usual smooth structure. Define
\[
  1\cdot x=x,\qquad s\cdot x=-x,\qquad
  \omega_x(u,v)=u_0v_1-u_1v_0.
\]
Thus $\omega=dx_0\wedge dx_1$. This is an alternating continuous
bilinear form varying smoothly with $x$. If $\omega_x(u,v)=0$ for every
$v$, testing $v=(0,1)$ and $v=(1,0)$ gives $u_0=u_1=0$, so it is
nondegenerate. In the global vector-space chart the exterior derivative
on constant vectors is
\[
 (d\omega)_x(u,v,w)=D\omega_x[u](v,w)
  -D\omega_x[v](u,w)+D\omega_x[w](u,v).
\]
The form is constant, so its actual Fr\'echet derivative vanishes and
$d\omega=0$. These are the usual symplectic-form conditions
\cite[\S1.3]{cannas}.

The action is jointly smooth: on each open component $\{g\}\times V$
it is the linear map $x\mapsto\varepsilon(g)x$, where
$\varepsilon(1)=1$ and $\varepsilon(s)=-1$. Its derivative is that same
linear map, and
\[
 \omega_{g\cdot x}(D(g\cdot)_x u,D(g\cdot)_x v)
  =\varepsilon(g)^2\omega_x(u,v)=\omega_x(u,v).
\]
It therefore acts by symplectomorphisms.

\paragraph{Lie algebra, coadjoint action and moment map.}
The actual tangent Lie algebra $\g=T_1G$ is zero-dimensional. It is a
zero vector space, so its dual $\g^*$ is also a singleton. The adjoint
map is defined as the derivative at $1$ of $h\mapsto ghg^{-1}$; the
coadjoint action is its inverse dual. Consequently the quotient of
$\g^*$ by this actual action has exactly one element.

Define $\mu:V\to\g^*$ by $\mu(x)=0$. It is smooth and equivariant.
For $\xi\in\g$, the fundamental vector field is the genuine orbit-map
derivative
\[
 \xi_V(x)=D_1(g\mapsto g\cdot x)[\xi].
\]
Since $\xi=0$, linearity of this derivative gives $\xi_V(x)=0$.
Each component $\mu^\xi(x)=\langle\mu(x),\xi\rangle$ is constant, hence
\[
 d\mu^\xi_x(v)=0=\omega_x(\xi_V(x),v).
\]
Together with equivariance, this is precisely the moment-map identity
\cite[Definition~22.1]{cannas}. In particular, $(V,\omega,G,\mu)$ is a
Hamiltonian $G$-space, and its zero level is all of $V$.

\paragraph{The two occurring types in the quotient.}
The reduced orbit space is the actual quotient
\[
 M_0=\mu^{-1}(0)/G=V/G.
\]
Its isotropy type at the orbit of $x$ is the conjugacy class of the
stabilizer $G_x=\{g:g\cdot x=x\}$. These are the standard orbit-type
indices \cite[Example~1.2 and Theorem~2.1]{sl}.
Direct calculation gives
\[
 G_0=G,\qquad G_x=\{1\}\quad(x\ne0),
\]
because $-x=x$ over $\R$ implies $x=0$. Both cases occur, for instance
at $0$ and $(1,0)$. The action is therefore nonfree. The two subgroups
are distinct and cannot be conjugate; indeed $G$ is abelian, so
conjugation fixes every subgroup. Isotropy type is invariant along
orbits and therefore defines a function on $M_0$. Its image is exactly
the two-element set $\{[G],[\{1\}]\}$. Thus
\[
 \#\{\text{occurring isotropy types in }M_0\}=2
 \ne1=\#(\g^*/G).
\]
The same mismatch holds if coadjoint orbits are restricted to the image
of $\mu$, since that image is already all of the singleton $\g^*$.

\paragraph{Formal scope.}
The Lean project uses the native \texttt{GroupLieAlgebra}, actual
\texttt{mfderiv} for the adjoint and fundamental fields, and the
continuous dual with its explicit linear equivalence to the algebraic
dual. The plane's canonical tangent spaces carry the displayed
bilinear form. Smoothness, nondegeneracy, the coordinate exterior
derivative, and the moment identity are proved from these definitions.
The zero level is a subtype, its reduction is the native
\texttt{MulAction.orbitRel} quotient, and its isotropy types come from
native \texttt{MulAction.stabilizer} subgroups modulo genuine
conjugacy. Finiteness and nonemptiness are proved before the two
cardinalities are compared. No cardinality is stipulated as a definition.

The result concerns the source's final count clause. It does not
contradict the classical orbit-type decomposition theorem, nor make a
claim about versions with extra connected-group or proper-moment-map
hypotheses. The full bilingual source and a pinned Lean 4.33.0 project
with Mathlib revision
\begin{center}\small
\texttt{db584cd6d46c92f209a44c0f1c829460d327499d}
\end{center}
are included, together with source-bound verification receipts.

\begin{thebibliography}{2}
\bibitem{cannas}
A.~Cannas da Silva, \emph{Lectures on Symplectic Geometry},
Lecture Notes in Mathematics 1764, Springer, 2001;
\href{https://people.math.ethz.ch/~acannas/Papers/lsg.pdf}{author version,
revised January 2006}, pp.~6, 133--135.
\bibitem{sl}
R.~Sjamaar and E.~Lerman, \emph{Stratified symplectic spaces and reduction},
Annals of Mathematics (2) 134 (1991), 375--422;
\href{https://people.maths.ox.ac.uk/beem/papers/symplectic_stratified_spaces.pdf}{author preprint},
Example~1.2 (p.~5), Theorem~2.1 (p.~9).
\end{thebibliography}
\end{document}
